% =========================================================================
%  Electromagnetismo I — Tema 2: Desarrollo multipolar
%  Estilo: artículo científico minimalista, idéntico al de los temas 0 y 1
%  (amsart, Palatino, versalitas, enunciados numerados por sección), con la
%  paleta cambiada a magentas y rojizos claros sobre grises sobrios. La
%  notación vectorial es la de «Cálculo integral multivariable» (L. Navas) y
%  la de los temas anteriores; las expresiones son las del formulario.
% =========================================================================

% --- IDIOMA Y CODIFICACIÓN ------------------------------------------------
\usepackage[utf8]{inputenc}
\usepackage[T1]{fontenc}
\usepackage[spanish,es-tabla,es-nodecimaldot,es-noshorthands]{babel}

% --- TIPOGRAFÍA (Palatino, como en los temas 0 y 1) -------------------------
\usepackage{mathpazo}
\linespread{1.05}
\emergencystretch=2em
\usepackage[final]{microtype}

% --- MATEMÁTICAS -------------------------------------------------------------
\usepackage{amsmath, amssymb, mathtools}
\usepackage{bm}
\usepackage{mathrsfs}
\usepackage{esint}

% --- PÁGINA, COLORES Y GRÁFICOS --------------------------------------------
\usepackage[a4paper,margin=2.6cm]{geometry}
\usepackage{graphicx}
\usepackage{xcolor}
% Paleta base del Tema 2: magentas sobrios y rojizos claros
\definecolor{magprof}{HTML}{5B1A45}     % magenta muy oscuro (títulos, enunciados)
\definecolor{magacero}{HTML}{A0306E}    % magenta medio (enlaces, banda de portada)
\definecolor{magclaro}{HTML}{E2A6C4}    % magenta claro (marcos)
\definecolor{magfondo}{HTML}{FDF3F8}    % fondo tenue de definiciones
\definecolor{rosaclaro}{HTML}{F6D5DE}   % rosa pálido (rellenos)
\definecolor{rojoclaro}{HTML}{D9645A}   % rojizo claro (acento)
\definecolor{grispizarra}{HTML}{6B7280} % gris
\definecolor{grisclaro}{HTML}{D4D8DE}   % gris claro
\colorlet{acento}{rojoclaro}
% Paleta ampliada para campos escalares y vectoriales (magnitud creciente:
% índigo -> púrpura -> carmín -> naranja -> ámbar), la misma de los temas 0 y 1
\definecolor{cm0}{RGB}{38,30,112}
\definecolor{cm1}{RGB}{112,36,140}
\definecolor{cm2}{RGB}{190,52,118}
\definecolor{cm3}{RGB}{236,104,66}
\definecolor{cm4}{RGB}{252,184,42}
% Colores de apoyo de los diagramas geométricos (tonos sobrios)
\definecolor{rojoteo}{HTML}{C2415B}     % vectores destacados, resultados (carmín suave)
\definecolor{verdeteo}{HTML}{2F6F68}    % versores y normales (verde azulado, contraste)
\definecolor{pizarra}{HTML}{3F4C6B}     % cargas negativas, ejes secundarios
\colorlet{ocreteo}{rosaclaro}           % superficies y volúmenes (relleno)
% Colores de las figuras clonadas de Mathematica (se conservan los originales)
\definecolor{mmafondo}{HTML}{E6EEF8}    % fondo azul pálido de StreamPlot
\definecolor{mmamarco}{HTML}{5B7DB8}    % marco de la región
\definecolor{mmalinea}{HTML}{1E1E9A}    % líneas de campo
\definecolor{mmaequi}{HTML}{4A4A4A}     % equipotenciales superpuestas
\definecolor{mmarojo}{HTML}{C00000}     % vector dipolar y cargas
\definecolor{mmaazul}{HTML}{3E7CB1}     % serie 1 (Mathematica 97)
\definecolor{mmanaranja}{HTML}{E1A027}  % serie 2
\definecolor{mmaverde}{HTML}{8FA93C}    % serie 3
\definecolor{mmarojizo}{HTML}{E0644A}   % serie 4
\definecolor{mmavioleta}{HTML}{8A7DB8}  % serie 5

\usepackage{tikz}
\usepackage{tikz-3dplot}
\usetikzlibrary{babel, arrows.meta, calc, positioning, decorations.pathmorphing,
                decorations.markings, patterns, 3d, intersections, angles, quotes,
                fadings, shapes.geometric, spy}
\usepackage{pgfplots}
\pgfplotsset{compat=1.17}
\usepgfplotslibrary{colormaps, patchplots, fillbetween}
\pgfplotsset{colormap={campo}{
  rgb255=(38,30,112) rgb255=(112,36,140) rgb255=(190,52,118)
  rgb255=(236,104,66) rgb255=(252,184,42) rgb255=(255,226,110)}}
\pgfplotsset{colormap={magenta}{
  rgb255=(253,243,248) rgb255=(226,166,196) rgb255=(160,48,110) rgb255=(91,26,69)}}
\usepackage{booktabs}
\usepackage{needspace}
\usepackage{array}
\usepackage{multirow}
\usepackage{enumitem}
\setlist[itemize]{leftmargin=1.6em, itemsep=2pt, topsep=4pt}
\setlist[enumerate]{leftmargin=1.8em, itemsep=2pt, topsep=4pt}

% --- ENLACES Y REFERENCIAS CRUZADAS ----------------------------------------
\usepackage[colorlinks=true, linkcolor=magacero, citecolor=magacero,
            urlcolor=magacero, plainpages=false, pdfusetitle=false,
            pdftitle={Tema 2. Desarrollo multipolar},
            pdfauthor={Daniel Martín Guerrero}]{hyperref}
\usepackage[spanish,nameinlink]{cleveref}

% --- BIBLIOGRAFÍA -----------------------------------------------------------
\usepackage{csquotes}
\usepackage[backend=biber, style=numeric, sorting=none, giveninits=true,
            maxbibnames=6, minbibnames=3, doi=true, url=true,
            isbn=false]{biblatex}
\addbibresource{referencias.bib}
\DeclareFieldFormat[article,techreport,unpublished,book,inproceedings]{title}{#1\isdot}
\DeclareFieldFormat{url}{\url{#1}}
\renewcommand*{\bibfont}{\small}
\setlength{\bibitemsep}{3pt}

% --- CABECERAS ---------------------------------------------------------------
\pagestyle{headings}
\makeatletter
\def\ps@headings{%
  \let\@mkboth\markboth
  \def\@evenhead{\small\scshape\normalfont\thepage\hfil Electromagnetismo I}%
  \def\@oddhead{\small\scshape\normalfont Tema 2. Desarrollo multipolar\hfil\thepage}%
  \let\@evenfoot\@empty \let\@oddfoot\@empty}
\makeatother
\numberwithin{equation}{section}
% Subsecciones como títulos en bloque (no «run-in»), con aire antes y después
\usepackage{titlesec}
\titleformat{\subsection}[block]{\normalfont\bfseries\color{magprof}}{\thesubsection.}{0.6em}{}
\titlespacing*{\subsection}{0pt}{3ex plus 1ex minus .5ex}{1.6ex plus .4ex}
\setcounter{tocdepth}{2}
\makeatletter
\def\@@and{y}
\renewcommand{\keywordsname}{Palabras clave}
% Índice con sangrías fijas: el número de sección nunca se monta sobre el título
\def\l@section{\@tocline{1}{5pt plus 1pt}{0pt}{1.9em}{\bfseries\color{magprof}}}
\def\l@subsection{\@tocline{2}{0pt}{1.9em}{2.7em}{}}
\renewcommand{\tocsection}[3]{\indentlabel{\@ifnotempty{#2}{\makebox[1.9em][l]{#2.}}}#3}
\renewcommand{\tocsubsection}[3]{\indentlabel{\@ifnotempty{#2}{\makebox[2.7em][l]{#2.}}}#3}
\makeatother

% =========================================================================
%               NOTACIÓN Y MACROS DE USO FRECUENTE
% =========================================================================
% --- Diferencial, derivadas parciales y números -----------------------------
\newcommand{\dif}{\mathrm{d}}
\newcommand{\pd}[2]{\frac{\partial #1}{\partial #2}}
\newcommand{\pdd}[2]{\frac{\partial^{2} #1}{\partial #2^{2}}}
\newcommand{\pdm}[3]{\frac{\partial^{2} #1}{\partial #2\,\partial #3}}
\newcommand{\R}{\mathbb{R}}
\newcommand{\N}{\mathbb{N}}
\newcommand{\Z}{\mathbb{Z}}
\newcommand{\eul}{\mathrm{e}}

% --- Vectores, versores y operador nabla ---------------------------------
\newcommand{\vc}[1]{\vec{#1}}                       % vector con flecha
\newcommand{\uh}[1]{\hat{u}_{#1}}                   % versor û_#1
\newcommand{\ux}{\uh{x}}\newcommand{\uy}{\uh{y}}\newcommand{\uz}{\uh{z}}
\newcommand{\ur}{\uh{r}}\newcommand{\uth}{\uh{\theta}}\newcommand{\uph}{\uh{\varphi}}
\newcommand{\urho}{\uh{\rho}}\newcommand{\uR}{\uh{R}}\newcommand{\un}{\hat{n}}
\newcommand{\nab}{\vec{\nabla}}                     % ∇⃗
\newcommand{\nabp}{\vec{\nabla}'}                   % ∇⃗′ (coordenadas fuente)
\newcommand{\Lap}{\nabla^{2}}                       % laplaciano
\newcommand{\Lapp}{\nabla'^{2}}                     % laplaciano en coordenadas fuente
\newcommand{\norm}[1]{\lVert #1\rVert}
\newcommand{\abs}[1]{\lvert #1\rvert}
\newcommand{\prodesc}[2]{\langle #1,#2\rangle}
\newcommand{\rp}{\vc{r}\,'}                         % r⃗′ punto fuente
\newcommand{\dl}{\dif\vc{\ell}}                     % elemento de línea dirigido
\newcommand{\dS}{\dif\vc{S}}                        % elemento de área dirigido
\newcommand{\dsc}{\dif S}                           % elemento de área escalar
\newcommand{\dV}{\dif V}                            % elemento de volumen
\newcommand{\defeq}{\mathrel{\overset{\mathrm{def}}{=}}}
\newcommand{\bordeS}{\partial S}
\newcommand{\bordeV}{\partial V}
\DeclareMathOperator{\grad}{grad}
\DeclareMathOperator{\dive}{div}
\DeclareMathOperator{\rot}{rot}
\DeclareMathOperator{\sen}{sen}
\DeclareMathOperator{\senh}{senh}
\DeclareMathOperator{\sgn}{sgn}
\DeclareMathOperator{\Tr}{Tr}
\renewcommand{\sin}{\sen}                            % «sen» en todo el texto
% Tensor de Levi-Civita y delta de Kronecker
\newcommand{\eps}{\varepsilon}
\newcommand{\dpart}[1]{\partial_{#1}}

% --- Notación de electrostática (Tema 1) -------------------------------------
\newcommand{\eo}{\varepsilon_{0}}                    % permitividad del vacío
\newcommand{\qp}{q\,'}                              % carga fuente q'
\newcommand{\pot}{\phi}                              % potencial electrostático
\newcommand{\Qenc}{Q_{\mathrm{enc}}}                 % carga encerrada
\newcommand{\Ke}{\frac{1}{4\pi\eo}}                  % constante de Coulomb 1/(4 pi eps0)
\newcommand{\vE}{\vc{E}}
\newcommand{\vF}{\vc{F}}
\newcommand{\dlp}{\dif \ell\,'}                     % elemento de longitud fuente
\newcommand{\dSp}{\dif S\,'}                        % elemento de superficie fuente
\newcommand{\dVp}{\dif V\,'}                        % elemento de volumen fuente
\newcommand{\ent}[1]{\bigl\lvert #1\bigr\rvert}    % módulo grande

% --- Notación específica del Tema 2 (desarrollo multipolar) ------------------
\newcommand{\vp}{\vc{p}}                             % momento dipolar
\newcommand{\vl}{\vc{\ell}}                          % separación del dipolo (de -q a +q)
\newcommand{\vP}{\vc{P}}                             % polarización
\newcommand{\rhoP}{\rho_{P}}                         % densidad volúmica de polarización
\newcommand{\sigP}{\sigma_{P}}                       % densidad superficial de polarización
\newcommand{\Pl}[1]{P_{#1}}                          % polinomio de Legendre P_ℓ
\newcommand{\cosp}{\cos\theta'}                      % coseno del ángulo entre r⃗ y r⃗′
\newcommand{\Qt}{\mathsf{Q}}                         % tensor cuadrupolar (matriz)
\newcommand{\Qjk}{Q_{jk}}                            % componentes del tensor cuadrupolar
\newcommand{\cd}[1]{\ell_{#1}}                       % cosenos directores ℓ_x, ℓ_y, ℓ_z
\newcommand{\phimul}[1]{\phi_{#1}}                   % término ℓ-ésimo del desarrollo
\newcommand{\OOp}{\overrightarrow{OO'}}              % traslación del origen
\newcommand{\delt}[2]{\delta_{#1#2}}                 % delta de Kronecker
\newcommand{\ord}{\mathcal{O}}                       % símbolo de orden de Landau

% Estilos de cargas puntuales para los diagramas
\tikzset{
  cargapos/.style={circle, fill=rojoteo, draw=magprof, line width=0.4pt, inner sep=0pt,
                   minimum size=9pt, text=white, font=\scriptsize\bfseries},
  carganeg/.style={circle, fill=pizarra, draw=pizarra!70!black, line width=0.4pt, inner sep=0pt,
                   minimum size=9pt, text=white, font=\scriptsize\bfseries},
  cargapru/.style={circle, fill=white, draw=magprof, line width=0.6pt, inner sep=0pt,
                   minimum size=9pt, text=magprof, font=\scriptsize\bfseries},
}

% =========================================================================
%        ENTORNOS FORMALES (enunciados numerados, estilo artículo)
% =========================================================================
\usepackage{amsthm}
\usepackage[most]{tcolorbox}

% Enunciados en cursiva (resultados generales)
\newtheoremstyle{cienplano}{10pt}{10pt}{\itshape}{}%
  {\bfseries\color{magprof}}{.}{0.6em}%
  {\thmname{#1}\thmnumber{ #2}\thmnote{\normalfont\upshape\ (#3)}}
% Definiciones, cálculos y ejemplos (cuerpo en redonda)
\newtheoremstyle{ciendef}{10pt}{10pt}{\normalfont}{}%
  {\bfseries\color{magprof}}{.}{0.6em}%
  {\thmname{#1}\thmnumber{ #2}\thmnote{\normalfont\upshape\ (#3)}}
% Observaciones y notas (encabezado en cursiva)
\newtheoremstyle{cienobs}{8pt}{8pt}{\normalfont}{}%
  {\itshape\color{magprof}}{.}{0.6em}%
  {\thmname{#1}\thmnumber{ #2}\thmnote{\normalfont\upshape\ (#3)}}

\theoremstyle{cienplano}
\newtheorem{proposicion}{Proposición}[section]
\newtheorem{teorema}[proposicion]{Teorema}
\newtheorem{lema}[proposicion]{Lema}
\newtheorem{corolario}[proposicion]{Corolario}
\theoremstyle{ciendef}
\newtheorem{notacion}[proposicion]{Notación}
\newtheorem{ejemplo}[proposicion]{Ejemplo}
\newtheorem{ejercicio}[proposicion]{Ejercicio}
% Definiciones: único entorno enmarcado (como en los temas 0 y 1): caja fina con
% fondo tenue y título sobre el borde. Comparte contador con los enunciados.
\newcounter{definicion}
\renewcommand{\thedefinicion}{\theproposicion}
\newtcolorbox{definicionbox}[1]{%
  enhanced, breakable, sharp corners,
  colback=magfondo, colframe=magclaro, boxrule=0.4pt,
  left=9pt, right=9pt, top=5pt, bottom=5pt,
  before skip=12pt, after skip=12pt,
  attach boxed title to top left={yshift*=-\tcboxedtitleheight/2, xshift=8pt},
  boxed title style={colback=white, colframe=white, sharp corners, boxrule=0pt,
                     left=3pt, right=3pt, top=1pt, bottom=1pt},
  fonttitle=\bfseries\color{magprof}, coltitle=magprof,
  title={Definición~\thedefinicion\ifx\relax#1\relax\else\ (#1)\fi}}
\newenvironment{definicion}[1][]
  {\par\stepcounter{proposicion}\refstepcounter{definicion}\begin{definicionbox}{#1}}
  {\end{definicionbox}}

\theoremstyle{cienobs}
\newtheorem{observacion}[proposicion]{Observación}

% Resolución de ejercicios: misma maquinaria que «proof», con otro rótulo
\renewcommand{\qedsymbol}{\ensuremath{\blacksquare}}
\newenvironment{resolucion}[1][Resolución]
  {\begin{proof}[\itshape\color{magprof}#1]}{\end{proof}}

\crefname{proposicion}{proposición}{proposiciones}
\Crefname{proposicion}{Proposición}{Proposiciones}
\crefname{teorema}{teorema}{teoremas}
\Crefname{teorema}{Teorema}{Teoremas}
\crefname{lema}{lema}{lemas}
\Crefname{lema}{Lema}{Lemas}
\crefname{corolario}{corolario}{corolarios}
\Crefname{corolario}{Corolario}{Corolarios}
\crefname{definicion}{definición}{definiciones}
\Crefname{definicion}{Definición}{Definiciones}
\crefname{notacion}{notación}{notaciones}
\Crefname{notacion}{Notación}{Notaciones}
\crefname{ejemplo}{ejemplo}{ejemplos}
\Crefname{ejemplo}{Ejemplo}{Ejemplos}
\crefname{ejercicio}{ejercicio}{ejercicios}
\Crefname{ejercicio}{Ejercicio}{Ejercicios}
\crefname{observacion}{observación}{observaciones}
\Crefname{observacion}{Observación}{Observaciones}
\crefname{section}{sección}{secciones}
\Crefname{section}{Sección}{Secciones}
\crefname{subsection}{sección}{secciones}
\Crefname{subsection}{Sección}{Secciones}
\crefname{table}{tabla}{tablas}
\Crefname{table}{Tabla}{Tablas}
\crefname{figure}{figura}{figuras}
\Crefname{figure}{Figura}{Figuras}
\crefname{equation}{ecuación}{ecuaciones}
\Crefname{equation}{Ecuación}{Ecuaciones}
\crefformat{equation}{ecuación~(#2#1#3)}
\Crefformat{equation}{Ecuación~(#2#1#3)}
\crefrangeformat{equation}{ecuaciones~(#3#1#4)--(#5#2#6)}
\crefmultiformat{equation}{ecuaciones~(#2#1#3)}{ y~(#2#1#3)}{, (#2#1#3)}{ y~(#2#1#3)}

% Referencia a un panel de figura: «figura 5\,(b)», enlazada en bloque
\newcommand{\figpanel}[2]{\hyperref[#1]{figura~\ref*{#1}\,\textup{(#2)}}}
\newcommand{\Figpanel}[2]{\hyperref[#1]{Figura~\ref*{#1}\,\textup{(#2)}}}

% Pies de figura y tabla, en la línea de los temas anteriores
\usepackage[font=small, labelfont={sc,color=magprof}, labelsep=period,
            margin=0.6cm, justification=justified]{caption}

% =========================================================================
%        ESTILOS GRÁFICOS COMUNES (TikZ / pgfplots)
% =========================================================================
% =========================================================================
%  Estilos y macros gráficas comunes del Tema 2
%  - \mapcolor{t}: escala cromática «campo» (t en [0,1]); define el color mapc
%  - Estilos TikZ heredados de los temas 0 y 1 (ejes, vectores, versores...)
%  - \panelMMA: marco de las figuras clonadas de Mathematica (StreamPlot,
%    ContourPlot y VectorPlot) que dibuja los datos de figuras/datos/*.tex
% =========================================================================

% --- Escala cromática (5 tramos entre cm0 y cm4) ------------------------
\newcommand{\mapcolor}[1]{%
  \pgfmathsetmacro{\mt}{min(max(#1,0),0.9999)*4}%
  \pgfmathtruncatemacro{\mi}{floor(\mt)}%
  \pgfmathtruncatemacro{\mj}{\mi+1}%
  \pgfmathsetmacro{\mf}{100*(\mt-\mi)}%
  \colorlet{mapc}{cm\mj!\mf!cm\mi}}

% --- Estilos TikZ generales ------------------------------------------------
\tikzset{
  >=Stealth,
  eje/.style={->, line width=0.6pt, grispizarra},
  ejeprof/.style={->, line width=0.7pt, magprof},
  vect/.style={->, line width=1pt, magprof},
  vectrojo/.style={->, line width=1pt, rojoteo},
  vectverde/.style={->, line width=1pt, verdeteo},
  fino/.style={line width=0.4pt, grispizarra},
  punteado/.style={densely dotted, line width=0.5pt, grispizarra},
  etiq/.style={font=\small, text=black!85},
  cuadro/.style={draw=grispizarra!60, line width=0.4pt, fill=white, inner sep=2pt},
  trayecto/.style={line width=1.1pt, rojoteo},
  orienta/.style={decoration={markings, mark=at position 0.55 with {\arrow{Stealth[length=5pt]}}}, postaction={decorate}},
  % líneas de campo y equipotenciales en la paleta del documento
  lcampo/.style={magacero, line width=0.55pt, postaction={decorate},
                 decoration={markings, mark=at position 0.5 with {\arrow{Stealth[length=3.6pt]}}}},
  lcampoinv/.style={magacero, line width=0.55pt, postaction={decorate},
                 decoration={markings, mark=at position 0.5 with {\arrowreversed{Stealth[length=3.6pt]}}}},
  lequi/.style={pizarra!70, line width=0.45pt, densely dashed},
}

% --- Estilo base de ejes pgfplots (como en los temas anteriores) -----------
\pgfplotsset{
  tema/.style={
    tick label style={font=\scriptsize}, label style={font=\small},
    axis lines=left, axis line style={grispizarra},
    grid=major, grid style={grisclaro, line width=0.4pt},
    every tick/.style={grispizarra},
    legend style={font=\scriptsize, draw=none, fill=white, fill opacity=0.85,
                  text opacity=1},
    /pgf/number format/use comma,
  },
  % Ejes con marco completo, como los «Plot» de Mathematica de las figuras del tema
  marcoMMA/.style={
    tick label style={font=\scriptsize}, label style={font=\small},
    axis line style={grispizarra!80, line width=0.5pt},
    every tick/.style={grispizarra}, tick align=inside,
    grid=major, grid style={grisclaro!70, line width=0.3pt},
    /pgf/number format/use comma,
    legend style={font=\scriptsize, draw=none, fill=white, fill opacity=0.85,
                  text opacity=1},
  },
  mapa/.style={
    tick label style={font=\scriptsize}, label style={font=\small},
    axis line style={grispizarra}, every tick/.style={grispizarra},
    /pgf/number format/use comma,
    colormap name=campo,
  },
}

% =========================================================================
%  Paneles «tipo Mathematica» (ventana [-L,L]^2 en unidades de \ell)
%  Uso dentro de un tikzpicture con x=y=<escala>:
%    \panelMMA{L}{fondo}{radio del disco central}{contenido}
%  El contenido se recorta a la ventana; luego se dibujan el disco blanco
%  (región excluida, donde están las fuentes) y el marco con marcas.
% =========================================================================
% Estilos de los objetos que escribe scripts/generar_datos.py
\tikzset{
  mmaLC/.style={mmalinea, line width=0.42pt},
  mmaEQ/.style={mmaequi, line width=0.28pt},
  mmaEQsolo/.style={black, line width=0.4pt},
  mmaFL/.style={-{Stealth[length=3.3pt, width=2.7pt]}, mmalinea, line width=0.42pt},
}
% Macros por defecto (cada figura puede redefinirlas antes de \input)
\newcommand{\EQ}[1]{\draw[mmaEQ] plot coordinates {#1};}
\newcommand{\LC}[1]{\draw[mmaLC] plot coordinates {#1};}
\newcommand{\FL}[3]{\draw[mmaFL] (#1,#2) -- ++(#3:0.02);}
\newcommand{\VF}[5]{%
  \mapcolor{#5}%
  \pgfmathsetmacro{\hl}{2.0+2.4*#5}\pgfmathsetmacro{\lw}{0.32+0.45*#5}%
  \draw[-{Stealth[length=\hl pt,width=0.85*\hl pt]}, line width=\lw pt, mapc]
     ({#1-0.5*#3},{#2-0.5*#4}) -- ++(#3,#4);}

% Marco con marcas en -L, -L/2, 0, L/2, L y rótulos de los ejes
% \marcoMMA{L}{paso de marcas}{rótulo x}{rótulo y}
\colorlet{marcoc}{mmamarco}
\newcommand{\marcoMMA}[4]{%
  \draw[marcoc, line width=0.55pt] (-#1,-#1) rectangle (#1,#1);
  \foreach \t in {-#1,-#2,0,#2,#1}{%
    \draw[grispizarra, line width=0.35pt] (\t,-#1) -- ++(0,0.025*#1*2);
    \draw[grispizarra, line width=0.35pt] (-#1,\t) -- ++(0.025*#1*2,0);
    \draw[grispizarra, line width=0.35pt] (\t,#1) -- ++(0,-0.025*#1*2);
    \draw[grispizarra, line width=0.35pt] (#1,\t) -- ++(-0.025*#1*2,0);
    \node[below=1pt, font=\tiny, inner sep=1pt] at (\t,-#1) {$\pgfmathprintnumber{\t}$};
    \node[left=1pt, font=\tiny, inner sep=1pt] at (-#1,\t) {$\pgfmathprintnumber{\t}$};}
  \node[font=\small, anchor=north] at (0,-#1) [yshift=-9pt] {#3};
  \node[font=\small, anchor=south, rotate=90] at (-#1,0) [xshift=0pt, yshift=12pt] {#4};}

% \panelMMA{L}{color de fondo}{r0}{contenido}
\newcommand{\panelMMA}[4]{%
  \fill[#2] (-#1,-#1) rectangle (#1,#1);
  \begin{scope}\clip (-#1,-#1) rectangle (#1,#1); #4 \end{scope}
  \ifdim #3pt>0pt
    \fill[white] (0,0) circle (#3);
    \draw[mmamarco!80, line width=0.4pt] (0,0) circle (#3);
  \fi}

% Vector dipolar rojo en el centro de un panel (altura ±h)
\newcommand{\dipoloRojo}[1]{%
  \draw[-{Stealth[length=4pt,width=3.6pt]}, mmarojo, line width=1.3pt] (0,-#1) -- (0,#1);}

% =========================================================================
%  Primitivas 3D para tdplot (vista fija: tdplotsetmaincoords{70}{125}),
%  heredadas del Tema 1.
% =========================================================================
\def\VA{125}
% \tdCirc[opciones]{radio}{z}   (mitad frontal continua, trasera a trazos)
\newcommand{\tdCirc}[3][]{%
  \draw[#1, domain=\VA-180:\VA, samples=36, variable=\tdt]
    plot ({#2*cos(\tdt)},{#2*sin(\tdt)},{#3});
  \draw[#1, densely dashed, domain=\VA:\VA+180, samples=36, variable=\tdt]
    plot ({#2*cos(\tdt)},{#2*sin(\tdt)},{#3});}
% \tdEsfera{radio}{color}  (esfera sombreada con ecuador)
\newcommand{\tdEsfera}[2]{%
  \begin{scope}[tdplot_screen_coords]
    \shade[ball color=rosaclaro!80, opacity=0.45] (0,0) circle (#1);
    \draw[#2, line width=0.7pt] (0,0) circle (#1);
  \end{scope}
  \tdCirc[#2!60, line width=0.4pt]{#1}{0}}

% =========================================================================
%  Región de carga V' para los diagramas geométricos (secciones 3 y 4):
%  contorno suave, separado de los ejes, con sombreado radial tenue.
%  \regionV{desplazamiento}  (el contorno está centrado cerca de (2.25,2.7))
% =========================================================================
\def\contornoV{plot[smooth cycle, tension=0.75] coordinates
  {(0.95,2.5) (1.35,3.5) (2.45,3.95) (3.55,3.45) (3.6,2.25) (2.6,1.55) (1.45,1.65)}}
\newcommand{\regionV}{%
  \shade[inner color=rosaclaro!25, outer color=rosaclaro!90] \contornoV;
  \draw[magacero, line width=0.75pt] \contornoV;}
% Ejes cartesianos en perspectiva caballera con el origen marcado
% \ejesO{long z}{long y}{long x}
\newcommand{\ejesO}[3]{%
  \draw[eje] (0,0) -- (0,#1) node[above, etiq] {$z$};
  \draw[eje] (0,0) -- (#2,0) node[right, etiq] {$y$};
  \draw[eje] (0,0) -- ({-0.76*#3},{-0.65*#3}) node[below left=-2pt, etiq] {$x$};}
% Elemento de volumen (cubito en perspectiva) centrado en #1
\newcommand{\cubito}[1]{%
  \draw[magprof, line width=0.5pt, fill=white] ($(#1)+(-0.12,-0.12)$) rectangle ($(#1)+(0.12,0.12)$);
  \draw[magprof, line width=0.5pt, fill=rosaclaro!40] ($(#1)+(-0.12,0.12)$) -- ++(0.07,0.07) -- ++(0.24,0) -- ++(-0.07,-0.07) -- cycle;
  \draw[magprof, line width=0.5pt, fill=rosaclaro!70] ($(#1)+(0.12,0.12)$) -- ++(0.07,0.07) -- ++(0,-0.24) -- ++(-0.07,-0.07) -- cycle;}

